\subsection{Informed trading and informative prices}

\begin{qfdefinition}[Fundamental value, market value and noise: what informative prices are informative about]\label{def:fundamental-value}
\qffield{Statement} The \emph{market value} of an instrument is the price at which traders can buy or sell it. Its \emph{fundamental value} (intrinsic value) is the expected present value of all present and future benefits and costs of holding it: the value everyone would agree on if all knew everything known about the instrument, analysed it properly and perceived its benefits and costs alike \citep[chap.~``Informed Traders and Market Efficiency'', sec.~10.1, pp.~222--223]{alpha-harris}. The difference between the two is \emph{noise}; prices are informative when they are near fundamental values and completely informative when they equal them \citep[same sec., p.~223]{alpha-harris}. Writing $V_t$ for fundamental value and $P_t$ for the price, noise is $\epsilon_t=P_t-V_t$.

Fundamental values use only information available now, so they are the best estimates of perfect-foresight values, not those values; and their changes are unpredictable, since a predictable change would mean they did not yet reflect the information behind the prediction \citep[same sec., p.~223]{alpha-harris}. Relative to a price, information is of three kinds:
\begin{itemize}
\item \emph{public}: available to any trader, such as public news and past prices and volumes \citep[same chap., pp.~223 and~240]{alpha-harris};
\item \emph{private}: held only by some traders, whether learnt first or produced by their own research \citep[same chap., pp.~223 and~228]{alpha-harris};
\item \emph{stale}: already in the price, meaning it cannot be used to forecast future price changes; traders who act on it are \emph{pseudo-informed}, and in fact uninformed \citep[same chap., sec.~10.5.2.1, p.~229]{alpha-harris}.
\end{itemize}

\qffield{Notation} $V_t$ fundamental value, $P_t$ market price, $\epsilon_t=P_t-V_t$ noise; the sign is a convention, since Harris defines noise only as the difference.

\qffield{Conditions} Fundamental values are not observed, so whether a price is informative or noisy cannot easily be determined \citep[same chap., p.~223]{alpha-harris}. Prices become informative in two ways: when values are common knowledge no one trades at other prices; otherwise informed traders' trading moves prices towards their estimates, which is the more important route because values are rarely common knowledge \citep[same chap., sec.~10.6.1, pp.~236--237]{alpha-harris}. Information gets into the price when all traders understand its significance or when informed traders push the price towards their estimates \citep[same chap., p.~229]{alpha-harris}.

\qffield{Practice} Whether information is stale can rarely be checked directly. Harris's guide: information that is unique and that others could not have anticipated is probably not stale; information acted on slowly, widely known, cheap to obtain, or reasonably anticipated probably is. His example: Occidental Petroleum rose 10 percent the day after its long-ill chairman died and fell back the next day, because his death was expected and its date was not material \citep[same chap., p.~229 and box, p.~230]{alpha-harris}. Fischer Black would call stock prices informative if they lay between half and twice their fundamental values \citep[same chap., box, p.~223]{alpha-harris}.

\qffield{Pitfalls} Two conditioning sets are in use. Harris's fundamental value conditions on everything known to anyone, private information included. Hasbrouck calls the expectation of the payoff conditional on \emph{all public information} the fundamental value or efficient price \citep[sec.~3.2, p.~26]{s6-hasbrouck}: the $m_t$ of the \hyperref[def:volatility-components]{random-walk decomposition}. Harris's noise is therefore not Hasbrouck's pricing error $s_t=p_t-m_t$. A price equal to the public-information expectation has a zero pricing error but nonzero noise whenever private information has not yet reached it (Interpretation). Fundamental value is also not the perfect-foresight value implied by all payoff-relevant information, which Harris distinguishes from it. Private information need not be restricted: research can produce it lawfully from public data \citep[chap.~``Insider Trading'', p.~585]{alpha-harris}.

\qffield{See also} \hyperref[def:volatility-components]{fundamental and transitory volatility} (the statistical version, conditioned on public information); \hyperref[res:informed-trading-aggregation]{informed trading and aggregation} (how prices become informative); \hyperref[def:market-efficiency]{market efficiency}; \hyperref[def:inside-information]{inside information}; \hyperref[def:outside-spread]{outside spread} (value traders' estimates in use).

\end{qfdefinition}

\begin{qfresult}[Informed trading makes prices informative: the price as an aggregate of value estimates]\label{res:informed-trading-aggregation}
\qffield{Statement} Informed traders buy when the price is below their estimate of fundamental value and sell when it is above. Their buying pushes prices up and their selling pushes them down, so their trading moves prices towards their estimates \citep[chap.~``Informed Traders and Market Efficiency'', secs.~10.2--10.3, pp.~223--224]{alpha-harris}. Harris's algebraic illustration \citep[same chap., box, p.~225]{alpha-harris}: $N$ traders hold unbiased forecasts $f_i=V+e_i$ of the value $V$, each wants the position $D_i=a(f_i-P)$, and the security is in zero net supply. Market clearing, $\sum_ia(f_i-P)=0$, gives
\[ P=\frac1N\sum_{i=1}^Nf_i=V+e_M, \qquad e_M=\frac1N\sum_{i=1}^Ne_i. \]
With independent errors, $e_M\to0$ as $N$ grows, by the law of large numbers: prices are most informative when many informed traders research independently. Interpretation: if the errors are also independent of $V$ with common variance $\sigma^2$, then $\mathbb{E}[P\mid V]=V$ and $\operatorname{Var}(P-V)=\sigma^2/N$.

\qffield{Conditions} When forecasts differ in accuracy, the best combination is a weighted average that gives more weight to the more accurate. Well-informed traders tend to be wealthier and take larger positions, so prices approximate such an optimally weighted average: the market is a statistical calculator \citep[same chap., p.~224 and box, p.~226]{alpha-harris}. Interpretation: with trader-specific slopes $a_i$ the same clearing condition gives $P=\sum_ia_if_i/\sum_ia_i$, still unbiased when each $f_i$ is; with independent errors of variances $\sigma_i^2$, $\operatorname{Var}(P-V)=\sum_ia_i^2\sigma_i^2/(\sum_ia_i)^2$, minimised by $a_i\propto1/\sigma_i^2$ at $\bigl(\sum_i\sigma_i^{-2}\bigr)^{-1}$, below every $\sigma_i^2$. Many mistaken traders, or one large one, make prices less informative, until their losses drive them out \citep[same chap., p.~224]{alpha-harris}.

Informed traders do not trade to make prices informative: price impact is their cost, and they want prices to adjust only after they have built their positions \citep[same chap., p.~224]{alpha-harris}. They profit only from uninformed traders; competition among them pushes prices closer to values and lowers their profits until the least successful just cover their costs \citep[same chap., secs.~10.6.3 and~10.6.5, pp.~237--239]{alpha-harris}.

\qffield{Practice} Trade aggressively when the information will soon be common knowledge or many others will act on it, since the first to trade profit most and a flood of orders makes liquidity suppliers wary; trade slowly (stealth trading) when the advantage will last, so that others cannot infer that the trader is well informed \citep[same chap., sec.~10.4, pp.~225--226]{alpha-harris}. Information already incorporated in the price is stale, and trading on it does not pay \citep[same chap., p.~229]{alpha-harris}.

The most profitable traders have estimates that are both precise and \emph{orthogonal}, uncorrelated with other traders'; they face little competition for liquidity when they trade. Precision and orthogonality trade off: a precise but crowded estimate and an imprecise but orthogonal one can be equally profitable \citep[same chap., sec.~10.6.2.1, p.~237]{alpha-harris}.

\qffield{Pitfalls} Aggregation removes independent errors, not shared ones. Interpretation: if the errors have common variance $\sigma^2$ and pairwise correlation $\rho$ (which requires $\rho\ge-1/(N-1)$), then $\operatorname{Var}(e_M)=\tfrac1{N^2}[N\sigma^2+N(N-1)\rho\sigma^2]=\sigma^2[\rho+(1-\rho)/N]\to\rho\sigma^2$ for $\rho\ge0$, so a shared data source, model or stale signal leaves an error that no number of traders removes; Harris notes that traders who use the same information and methods have highly correlated estimates \citep[same chap., p.~237]{alpha-harris}. Harris also says the market's forecast error is smaller than the average individual error whenever the individual errors are not identical, and that the price estimates value better than any individual trader can \citep[same chap., box, p.~225]{alpha-harris}. Interpretation, by error measure: the signed errors average exactly, $e_M=\tfrac1N\sum_ie_i$, so ``smaller'' needs a loss.

For one realisation, squared error gives $e_M^2\le\tfrac1N\sum_ie_i^2$, since the gap is $\tfrac1N\sum_i(e_i-e_M)^2$, with equality only when all $e_i$ are equal, as Harris's condition says. Absolute error gives $|e_M|\le\tfrac1N\sum_i|e_i|$ by the triangle inequality, with equality whenever the errors share a sign, so $e=(1,2)$ is not identical yet gives $1.5$ on both sides. In expectation, with independent, zero-mean, nondegenerate errors, both inequalities are strict, since errors of opposite sign then occur with positive probability. With independent, unbiased errors, equal accuracy is sufficient for the equal-weighted average to beat every individual in expected squared error, but not necessary: with variances $1$ and $2$ the average has variance $0.75$, below both. It can also lose to the best individual: with variances $0.01$ and $1$ the average has variance $0.2525$, worse than the better trader's, while inverse-variance weights give $0.0099$.

The assumption is that the price is the average forecast; in the box this follows from linear demands with a common slope and zero net supply, an assumption Harris says simplifies the arithmetic without changing the qualitative results. Informed traders act on estimates of value, not on value itself, and lose when their estimates are poor \citep[same chap., p.~223]{alpha-harris}.

\qffield{See also} \hyperref[def:fundamental-value]{fundamental value}; \hyperref[def:market-efficiency]{market efficiency}; \hyperref[mod:kyle-single-auction]{Kyle single-auction model} and \hyperref[mod:glosten-milgrom]{Glosten--Milgrom model} (equilibrium models of informed trading); \hyperref[def:price-formation-process]{price formation process}; \hyperref[def:informed-trader-types]{informed-trader types}; \hyperref[def:order-anticipation]{order anticipation} (front runners accelerate informed traders' impact).

\end{qfresult}

\begin{qfdefinition}[Informed-trader types: value traders, news traders, information-oriented technical traders and arbitrageurs]\label{def:informed-trader-types}
\qffield{Statement} Harris classifies informed traders by what they estimate \citep[chap.~``Informed Traders and Market Efficiency'', sec.~10.5, pp.~226--235, Table~10-3, p.~235]{alpha-harris}:
\begin{itemize}
\item \emph{value traders} estimate total fundamental value from all available information; slow, because careful review guards against biased estimates;
\item \emph{news traders} estimate changes in value from new information, assuming the price already reflects everything else; fast on public news, slower on information they produce by their own research;
\item \emph{information-oriented technical traders} identify price patterns left by informed traders' systematic mistakes, such as under- or overreaction; fast;
\item \emph{arbitrageurs} estimate relative values of instruments that share valuation factors; fast.
\end{itemize}

\qffield{Conditions} Value and news traders profit in different circumstances. Price departs from value either because value changes and price lags, where news traders, first to learn the news, profit; or because uninformed traders push price away from value, where value traders profit \citep[chap.~``Value Traders'', sec.~16.1, pp.~338--339]{alpha-harris}. Their mistakes put them on opposite sides. News traders who misjudge the size of a change leave underreaction for value traders to trade on, or overreaction for them to trade against (Table~10-1, p.~230); pseudo-informed traders acting on stale news lose to value traders (box, p.~229). Value traders who mistake a price move driven by news for uninformed pressure lose to news traders \citep[chap.~``Informed Traders and Market Efficiency'', p.~239]{alpha-harris}.

\qffield{Practice} Organisation follows speed: large value traders have pyramid-shaped organisations with layers of review, news traders flat organisations that decide quickly \citep[same chap., pp.~227--229]{alpha-harris}. Technical traders who trade on the predictable impact of \emph{uninformed} traders are not informed traders: they act as dealers if they supply liquidity, or as order anticipators if they front-run \citep[same chap., p.~231]{alpha-harris}. Information-oriented technical strategies decay as informed traders learn from their mistakes \citep[same chap., p.~231]{alpha-harris}.

\emph{Liquidity roles}. Harris's Table~8-4 lists value traders, news traders, information-oriented technical traders and arbitrageurs as the informed traders \citep[chap.~``Why People Trade'', sec.~8.2.1, pp.~194--196]{alpha-harris}. As liquidity suppliers, value traders trade when uninformed demand has moved prices away from value and are ``the ultimate suppliers of liquidity'', supplying depth, while news traders trade on changes in value before prices adjust \citep[chap.~``Liquidity'', sec.~19.3.4, pp.~402--403]{alpha-harris}; arbitrageurs move liquidity between markets (\hyperref[def:arbitrage]{arbitrage}). The distinction, whether an order corrects stale prices, absorbs non-value flow or links related prices, maps onto news traders, value traders and arbitrageurs.

\qffield{Pitfalls} Technical traders do not in general exploit predictable errors by other traders. In Harris that holds for errors by informed traders; errors by uninformed traders belong to sentiment-oriented technical traders, who are order anticipators (\hyperref[def:order-anticipation]{order anticipation}). Dealer inventory pressure is not a separate source of value-trader profit: in Harris it is a channel of uninformed demand (see \hyperref[def:outside-spread]{outside spread}).

\qffield{See also} \hyperref[res:informed-trading-aggregation]{informed trading and aggregation}; \hyperref[def:outside-spread]{outside spread}; \hyperref[def:arbitrage]{arbitrage}; \hyperref[def:order-anticipation]{order anticipation}; \hyperref[def:fundamental-value]{fundamental value} (stale information).

\end{qfdefinition}

\begin{qfdefinition}[Winner's curse: what winning a common-value trade reveals, and how bids adjust]\label{def:winners-curse}
\qffield{Statement} Buyers who compete for an item of common but unknown value suffer the \emph{winner's curse} when they pay more than it is worth. The highest bidder tends to be the one who most overestimated the value, so bids near one's own estimate lose on average when they win. Buyers avoid it by lowering their bids to reflect what winning would tell them about their estimate; sellers whose offers are taken symmetrically raise their offers \citep[chap.~``Value Traders'', sec.~16.2.1.2, pp.~341--342]{alpha-harris}. Interpretation: let $\widehat V$ be the buyer's estimate, $c$ an all-in cost, and $\omega(b)=\mathbb{E}[\widehat V-V\mid\widehat V,\ \text{win at bid }b]$ the expected overestimate given one's own estimate and winning at bid $b$. The expected profit of winning is $\widehat V-\omega(b)-b-c$, so the highest bid that breaks even solves $b=\widehat V-\omega(b)-c$.

\qffield{Conditions} The required adjustment is larger \citep[same sec., p.~342]{alpha-harris}:
\begin{itemize}
\item the more bidders there are, because outbidding many implies a larger overestimate;
\item the harder the value is to estimate, because extreme estimates are then further from the truth.
\end{itemize}
Against market power: with only one or very few competitors a bidder may underbid in the hope of winning cheaply, so bids are more aggressive against several competitors than against one or two; this matters only with very few \citep[same sec., p.~342]{alpha-harris}. In a two-sided single-price auction the curse matters little, because traders do not trade at their own bids and offers and the clearing price is near the median estimate. In continuous markets they do, and aggregation is sequential, so traders who misestimate lose when they push prices \citep[same sec., pp.~342--343]{alpha-harris}.

\qffield{Practice} Oil companies overpaid for drilling leases until they allowed for the curse; their geologists were right on average over all tracts but optimistic about the tracts they won \citep[same chap., box, p.~342]{alpha-harris}. Value traders face the curse whenever they trade on a large apparent mispricing and widen their outside spreads in response; Harris attributes part of late-1990s Internet-stock volatility to prices bouncing within wide outside spreads \citep[same chap., p.~343 and box]{alpha-harris}. Do not compete against foolish bidders: one cannot profit bidding against strategies that ensure losses \citep[same chap., p.~342]{alpha-harris}.

\qffield{Pitfalls} The break-even bid $b=\widehat V-\omega(b)-c$ conditions on one's own estimate and on winning, which depends on the bid itself and, in equilibrium, on rivals' strategies, so it is a fixed-point condition, not a closed form. It also gives the highest acceptable bid, not the profit-maximising one, which shades further when competitors are few (Interpretation). Shared data make estimates move together, which reduces how much winning reveals, while independent noisy estimates increase selection risk (Interpretation; not in the source).

\qffield{See also} \hyperref[def:outside-spread]{outside spread} (the curse widens value traders' spreads); \hyperref[res:informed-trading-aggregation]{informed trading and aggregation} (correlated and independent estimates); \hyperref[mod:glosten-milgrom]{Glosten--Milgrom model} (adverse selection, the analogous inference a dealer makes from the side of a trade).

\end{qfdefinition}

\begin{qfdefinition}[Market efficiency: weak, semistrong and strong forms, and the definition with information costs]\label{def:market-efficiency}
\qffield{Statement} In an \emph{efficient market}, prices reflect all information that traders can acquire and profitably trade upon. The definition builds in the costs of acquiring information and of acting on it, and the impact of informed trading on prices \citep[chap.~``Informed Traders and Market Efficiency'', sec.~10.6.6, p.~240]{alpha-harris}. Harris sets it beside the three traditional forms \citep[same sec., p.~240]{alpha-harris}:
\begin{itemize}
\item \emph{weak form}: prices reflect all information in past prices, so no one can predict price changes from past prices alone;
\item \emph{semistrong form}: prices reflect all public information, which includes all public news and past prices and volumes in all securities and contracts;
\item \emph{strong form}: prices reflect all public and private information as soon as it is known; since informed traders could never profit, only markets in instruments whose values are common knowledge are strong-form efficient.
\end{itemize}

\qffield{Conditions} Timeless (concepts); the empirical summary as of 2002. Prices never fully reflect all the information traders could collect: information that is too expensive or too inconsequential is not acquired \citep[same sec., p.~240]{alpha-harris}. Harris's market paradox: if prices were fully informative, informed trading would not pay, so no one would make them informative. Prices and values change constantly, so gaps open and informed trading closes them: prices are informative but not always \citep[same chap., sec.~10.6.4, pp.~238--239]{alpha-harris}.

Lehalle and Laruelle make the same point with Grossman and Stiglitz (1980): market dynamics must contain enough inefficiency to reward those who improve the information in prices \citep[Introduction, p.~4]{alpha-laruelle}. How informative prices are depends on the costs of information and of trading, and on there being uninformed traders to lose to informed ones \citep[same chap., secs.~10.6.3 and~10.7, pp.~238 and~243--244]{alpha-harris}.

\qffield{Practice} Informative prices are not free. What uninformed traders lose to informed traders pays for much of the information in prices, so policies that frustrate informed traders can make prices less informative and markets more liquid; restrictions on insider trading are an example \citep[same chap., sec.~10.6.7, p.~241]{alpha-harris}.

The trade-off is not universal: publishing material fundamental information raises both liquidity and informativeness \citep[chap.~``Informed Traders and Market Efficiency'', sec.~10.6.8, pp.~241--242]{alpha-harris}. Trading halts before material releases protect uninformed traders, especially those who offer liquidity, at the cost of a few minutes' delay in price formation, a trade Harris judges worth making \citep[same chap., box, p.~241]{alpha-harris}. Most published studies find markets weak-form efficient, and semistrong-form efficient for public information that is easy to obtain and interpret \citep[same chap., p.~240]{alpha-harris}.

\emph{Liquidity and informative prices}. Liquid markets make informed trading cheaper, so informed traders can afford more research and trade on smaller pieces of information, and ``Prices therefore should be more informative in liquid markets than in illiquid markets''; informed traders profit only if uninformed traders are willing to lose to them \citep[chap.~``Good Markets'', sec.~9.5.4, pp.~215--216]{alpha-harris}. Visibility, update speed and uninformed flow therefore decide how much information enters prices, through this mechanism and that of \hyperref[res:informed-trading-aggregation]{informed trading and aggregation}.

\qffield{Pitfalls} Comparing market structures needs the change in benefits, not their totals, and the benefits of informative prices and of liquidity are both hard to measure \citep[same chap., sec.~10.6.7, p.~241]{alpha-harris}.

\qffield{See also} \hyperref[def:fundamental-value]{fundamental value}; \hyperref[res:informed-trading-aggregation]{informed trading and aggregation}; \hyperref[def:inside-information]{inside information} (the main example of the trade-off); \hyperref[res:spread-determinants]{spread determinants} (asymmetric information widens spreads); \hyperref[def:order-anticipation]{order anticipation} (front running of informed traders lowers informativeness in the long run).


\end{qfdefinition}

\begin{qfdefinition}[Inside information: material nonpublic information and who may not trade on it (United States, as summarised in 2002)]\label{def:inside-information}
\qffield{Statement} For economic analysis, Harris defines \emph{inside information} as material information about the value of a security that is not available to public traders, and \emph{material} information as information that would cause prices to change if it were widely known; \emph{insider trading} is trading on it \citep[chap.~``Insider Trading'', sec.~29.1, p.~584]{alpha-harris}. Inside information keeps its status however many people pass it on, and loses it only when it is publicly available, through a broadly distributed press release or a public filing \citep[same sec., p.~585]{alpha-harris}.

The legal definitions come from statutes, regulations and case law, are complex, and call for an attorney \citep[same sec., p.~584]{alpha-harris}. The SEC's summary, as quoted by Harris from its web page of January 2002, covers trading in breach of a fiduciary duty or other relationship of trust and confidence while in possession of material nonpublic information, and tipping, trading by those tipped, and misappropriation \citep[same sec., box, p.~585]{alpha-harris}. Those who may not trade on such information include:
\begin{itemize}
\item \emph{corporate insiders}: senior managers, directors, large shareholders, key employees, and anyone involved in corporate activities that will materially affect the firm's value \citep[sec.~29.1.1, p.~587]{alpha-harris};
\item employees of law, banking, brokerage and printing firms given the information to serve the corporation, and government employees who learn it at work \citep[sec.~29.1, box, p.~585]{alpha-harris};
\item \emph{tippees}: friends, family and business associates of insiders, and the friends of their friends \citep[same sec., pp.~584--585]{alpha-harris};
\item \emph{misappropriators}, who take confidential information from an employer or other source. Chiarella's conviction was reversed in 1980 because he owed no fiduciary duty to the firms whose stock he traded; Harris says that under the later law he would be convicted for misappropriation \citep[same sec., box, p.~586]{alpha-harris}. Winans was convicted under the misappropriation doctrine \citep[same sec., box, p.~587]{alpha-harris}.
\end{itemize}
Analysts whose private information comes only from analysing public information are not trading on inside information \citep[same sec., p.~585]{alpha-harris}. Harris's box ``Inside information, skilled speculation, or good advice?'' shows why the line is hard to draw from the trade alone: the same well-timed purchase before a takeover may rest on a tip, on independent analysis of public information, or on analysis that the trader then sold to the target's management; and a trader who had a tip will claim the second \citep[sec.~29.2, box, p.~590]{alpha-harris}.

\qffield{Conditions} Jurisdiction and date: United States law as Harris summarises it in 2002. Most countries prohibit insider trading, but only a few, mainly the United States, Canada and Great Britain, regularly and seriously enforce their laws \citep[chap.~``Insider Trading'', p.~584]{alpha-harris}. Regulation FD (2000) requires an issuer that discloses material nonpublic information to anyone to disclose it publicly \citep[same chap., box, p.~586]{alpha-harris}.

\qffield{Practice} Corporations restrict their employees' trading: many require corporate approval before a trade, may refuse it before a release, or allow trading only in a window after quarterly earnings \citep[same chap., sec.~29.1.1, p.~588]{alpha-harris}. Exchange surveillance looks for large price changes on large volume before significant news \citep[same chap., sec.~29.2, p.~588]{alpha-harris}. Insider trading is hard to prove: a sophisticated speculator can claim independent analysis of public data, and trading on information that is not revealed in a discrete event is essentially undetectable from prices and volumes \citep[same chap., pp.~589--590]{alpha-harris}. Like all informed trading, insider trading hurts liquidity suppliers and raises spreads for uninformed traders \citep[same chap., sec.~29.3.1.2, pp.~591--592]{alpha-harris}.

\emph{Controls on insiders' own trading}. Harris lists, for the United States in 2002, reporting of insider trades to the SEC, prearranged trading programmes that let insiders trade while knowing material information, return of short-swing profits on positions opened and closed within six months, restrictions on short positions in the employer's stock, corporate pre-approval of trades, and trading windows after quarterly earnings (sec.~29.1.1, pp.~587--588).

Current U.S.\ rules, as of the eCFR text for 14 October 2026: Rule 10b5-1 gives an affirmative defence for trades under a plan adopted before the trader was aware of the information; since its 2022 amendments (effective 27 February 2023) directors and officers must wait out a cooling-off period of the later of 90 days or two business days after the issuer's results for the quarter in which the plan was adopted are disclosed, up to 120 days, and other persons except the issuer 30 days \citep[paras.~(b) and (c)(1)(ii)(B)]{cfr-240-10b5-1} \citep[summary and dates, p.~80362]{sec-34-96492}; the same release added disclosure of plan adoptions and terminations and a Form~4 flag for plan trades (summary, p.~80362).

Form~4 must be filed ``before the end of the second business day following the day on which the subject transaction has been executed'' \citep[para.~(g)(1)]{cfr-240-16a-3}. Regulation~FD requires public disclosure ``simultaneously'' for intentional and ``promptly'' for non-intentional selective disclosures \citep[para.~(a)]{cfr-243-100}.

\qffield{Pitfalls} Harris's categories are legal categories summarised for economists, not a legal test; check the law in force. The reporting deadlines, prearranged-plan rules and penalties he gives \citep[same chap., pp.~587--591]{alpha-harris} are as of 2002; the reporting deadline and the plan rules here reflect the October 2026 rules above, while the penalties remain those of the 2002 account. The terms ``temporary insider'' and a ``relationship, benefit or confidence'' test for tipping do not appear in the source and need separate legal sources, as does the status of auditors, consultants and exchange staff, whom the SEC summary does not list. Whether a predictive signal may be used depends on data provenance, confidentiality duties, firm restricted lists, compliance approval and execution capacity.

The source supports confidentiality duties (misappropriation), employer approval of employees' own trades (p.~588), and the dependence of informed trading profits on liquidity \citep[chap.~``Informed Traders and Market Efficiency'', box, p.~236]{alpha-harris}. Firm restricted lists, compliance approval of signals, the provenance of purchased data, and the rights status of alternative data and channel checks are not discussed there. Two of Harris's statements are superseded: the ten-day reporting deadlines (p.~587) have been replaced by the two-business-day Form~4 deadline, and ``insiders do not have to disclose these plans to the public'' (p.~588) no longer holds after the 2022 amendments. Disclosure controls, ownership reports, prearranged plans, blackout windows, short-swing rules, surveillance and audit trails summarise the controls; blackout windows are Harris's trading windows seen from the other side.

\qffield{See also} \hyperref[def:fundamental-value]{fundamental value} (private and public information); \hyperref[def:market-efficiency]{market efficiency} (insider-trading restrictions in the information--liquidity trade-off); \hyperref[res:spread-determinants]{spread determinants} (enforced insider-trading rules as a secondary determinant); \hyperref[res:insider-trading-deterrence]{insider-trading profitability, detection and deterrence}.

\end{qfdefinition}

\begin{qfresult}[Insider-trading profitability, detection and deterrence, and the policy debate (United States, as described in 2002)]\label{res:insider-trading-deterrence}
\qffield{Statement} Harris's account of when trading on inside information pays and when it is caught \citep[secs.~29.2--29.3, pp.~588--595]{alpha-harris}:
\begin{itemize}
\item \emph{competition among insiders}: one insider with information that will not soon be released can trade slowly and extract the most; many insiders who learn it at once must race, which quickly reveals the information and lowers their total profit \citep[sec.~29.3.1.2, p.~592]{alpha-harris};
\item \emph{discreteness of release}: trading on information revealed in a discrete event can be traced back from it, whereas trading on information that firms do not reveal for months or years is essentially undetectable from prices and volumes \citep[sec.~29.2, p.~590]{alpha-harris}; restrictions are therefore effective mainly against information that will soon be public \citep[sec.~29.3.2.1, p.~595]{alpha-harris};
\item \emph{deterrence}: a rational criminal trades when the expected benefit exceeds the expected cost, the probability of being caught times the penalty, so a low probability of detection requires a large penalty \citep[sec.~29.2, p.~590]{alpha-harris}.
\end{itemize}
Interpretation: with expected gross gain $G$, execution and leakage costs $K$, a risk charge $R$, detection probability $q$ and penalty $F$, the trade is deterred when $G\le K+R+qF$. If the penalty were only disgorgement of the gain, $F=G$, the condition becomes $(1-q)G\le K+R$, which a small $q$ rarely satisfies; hence fines on top of disgorgement.

\qffield{Conditions} United States as Harris describes it in 2002. Penalties then included disgorgement, fines and, rarely, prison; they were heavier for traders closer to the source. He gives maximum criminal penalties of ten years and a million dollars, a maximum civil penalty of three times the profit gained or loss avoided, and an SEC bounty of up to 10 percent of civil penalties collected \citep[sec.~29.2, pp.~590--591 and box]{alpha-harris}. These figures are dated and unverified against current law. $R$ is a cost only for a risk-averse trader; a risk-neutral trader would not subtract it (Interpretation).

\qffield{Practice} Exchange surveillance identifies most cases. It looks for large price changes on large volume and asks the listed firm whether news explains them; trading before releases of significant unexpected news often indicates insiders or confederates \citep[sec.~29.2, p.~588]{alpha-harris}. Investigators then match a list of those who knew with a list of those who profited, in stocks, options or bonds, through links such as family, neighbours or classmates \citep[same sec., pp.~588--589 and box]{alpha-harris}. In the Colt Industries case, deep out-of-the-money options bought by friends and relatives of a law-firm employee were traced through the buyers' home addresses \citep[same sec., box, p.~589]{alpha-harris}.

\emph{The policy debate}. Harris's arguments for restricting insider trading: fairness and investor confidence (rebutted by the claim that firms can restrict it by contract, and answered by the difficulty of private enforcement), liquidity (insiders, like all informed traders, widen spreads and hurt limit-order traders), and corporate control (insiders who may trade share less information, may choose projects for their trading advantage, and may front-run their firms' own trades); against: informative prices (answered by the fact that enforcement works mainly against information soon to be public, so the efficiency gain is small), costs of enforcement and selective prosecution, and Manne's entrepreneurial incentives (answered by the need for positions to be taken before outcomes are known, and the general ban on insiders' short positions) \citep[sec.~29.3, pp.~591--597]{alpha-harris}.

His summary: the laws ``narrow bid/ask spreads and make prices less informative in the short run'' (pp.~597--598). Saying that a restriction improves welfare if and only if its liquidity and control gains exceed enforcement and delay costs, $L+C>E+D$, is an accounting identity that lists these terms; its content lies in their sizes, which Harris does not quantify. Enforcement detects abnormal trading before an event, lists who knew, matches traders to access, tests alternative explanations and then sanctions, as in the surveillance and investigation steps above.

\qffield{Pitfalls} These footprints are leads, not proof. Relationships between the two lists do not prove insider trading, and a sophisticated speculator will claim independent analysis of public data \citep[same sec., pp.~589--590]{alpha-harris}. Widened spreads, quote fading and dark-pool activity before news are not among Harris's signals, although wider spreads in the presence of insiders follow from adverse selection \citep[sec.~29.3.1.2, p.~591]{alpha-harris}. Market liquidity as a factor in insider profits is Harris's point for informed trading generally (chap.~``Informed Traders and Market Efficiency'', box, p.~236).

\qffield{See also} \hyperref[def:inside-information]{inside information} (definition and who is covered); \hyperref[res:informed-trading-aggregation]{informed trading and aggregation} (aggressive against stealth trading); \hyperref[def:market-efficiency]{market efficiency} (restrictions in the information--liquidity trade-off); \hyperref[def:bluffing-manipulation]{bluffing and manipulation} (the same proof problem).

\end{qfresult}

\begin{qfdefinition}[Price formation process: immediate impact and its decay]\label{def:price-formation-process}
\qffield{Statement} The \emph{price formation process} (PFP) is how trades and information become prices. Lehalle and Laruelle report a consensus that it consists of the impact of each trade combined with the variation of the asset's ``fair price'', after Bouchaud et al.\ (2002) \citep[Introduction, p.~27]{alpha-laruelle}. They split it into two phases \citep[same, Figure~8, pp.~27--29]{alpha-laruelle}:
\begin{enumerate}
\item \emph{immediate impact}: trades move the price through the balance of buyers and sellers, and of liquidity-providing and liquidity-consuming orders, together with moves of the spread;
\item \emph{decay}, or the \emph{resilience} of the book: liquidity returns, or does not, according to the distance between the best quotes and investors' views of the fair price.
\end{enumerate}
Pressure inconsistent with participants' consensus generates only temporary impact; pressure consistent with it is not pushed back, and the impact is permanent \citep[same, pp.~3 and~28]{alpha-laruelle}. The PFP is mainly driven by information \citep[same, p.~7]{alpha-laruelle}.

\qffield{Conditions} Timeless (concept); the regulatory context is European, 2007--2018. The more investors with views of the fair price take part, the faster impact decays. Market makers with inventory constraints contribute to the decay without any view of the fair price, but only if their inventory exceeds the turnover of a price jump \citep[Introduction, p.~28]{alpha-laruelle}. Post-trade transparency speeds the decay by warning traders that the last price is out of line with their views \citep[same, p.~28]{alpha-laruelle}. A design tension runs through it: the more information traders have, the better the PFP should be, but each trader fears information leakage, front-running and excessive impact \citep[same, pp.~19--20]{alpha-laruelle}.

\qffield{Practice} Market makers buy from early sellers and sell to later buyers, so that most temporary oscillations shrink to a spread related to the asset's risk \citep[Introduction, pp.~3--4]{alpha-laruelle}. Volatility is the random, unexpected component of the PFP, and the more efficient the microstructure, the lower intraday volatility \citep[chap.~2, sec.~2.5.1, pp.~171--172]{alpha-laruelle}. In a fragmented market, a consolidated tape organises how information is shared \citep[Introduction, p.~7]{alpha-laruelle}. The authors see high-frequency traders' growing role as a possible threat to the decay, since they need not hold a view of the fair price and have small inventory limits \citep[same, p.~28]{alpha-laruelle}; that question belongs to \nolinkurl{hft-market-design}.

\qffield{Pitfalls} The impact of a single large order depends on what others are doing. If many others buy at the same time, the impact of a large buy is permanent; if most sell, it can be almost invisible. Impact is measured only by averaging over many market configurations \citep[Introduction, p.~3]{alpha-laruelle}. ``Clean'' and ``noisy'' PFP, these terms, are not defined by the authors, although they speak of a ``less noisy PFP'' \citep[chap.~2, sec.~2.5.1, p.~177]{alpha-laruelle}. Interpretation: the noise is the transitory component of \hyperref[def:volatility-components]{price changes}.

\qffield{See also} \hyperref[res:informed-trading-aggregation]{informed trading and aggregation} (the information side); \hyperref[res:spread-component-identification]{component identification} (permanent and transitory price effects of trades); \hyperref[mod:kyle-single-auction]{Kyle single-auction model} (permanent impact from information); \hyperref[mod:grossman-miller]{Grossman--Miller model} (temporary concession for immediacy); \hyperref[def:intraday-overnight-returns]{returns by time scale}; \hyperref[def:price-impact]{price impact} (reaction and prediction impact, and the two readings of impact).

\qffield{Further reading} Bouchaud et al.\ (2002), cited for the decomposition.

\end{qfdefinition}
